% \section{Introduction}
\label{section:1}

The proliferation of visualizations during the {\smaller COVID-19} pandemic has underscored their double-edged potential: efficiently communicating critical public health information\,---\,as with the immediately-canonical ``Flatten the Curve'' chart (Fig.~\ref{fig:teaser})\,---\,while simultaneously excluding people with disabilities.
\say{\emph{For many people with various types of disabilities, graphics and the information conveyed in them is hard to read and understand,}} says software engineer Tyler Littlefield~\cite{ehrenkranz_vital_2020}, who built a popular text-based {\smaller COVID-19} statistics tracker after being deluged with inaccessible infographics~\cite{sutton_accessible_2020, littlefield_covid-19_2020}.
While natural language descriptions sometimes accompany visualizations in the form of chart captions or alt text (short for \say{alternative text}), these practices remain rare.
Technology educator and researcher Chancey Fleet notes that infographics and charts usually lack meaningful and detailed descriptions, leaving disabled people with \say{\emph{a feeling of uncertainty}} about the pandemic~\cite{ehrenkranz_vital_2020}.
For readers with visual disabilities (approximately 8.1 million in the United States and 253 million worldwide~\cite{ackland_world_2017}), inaccessible visualizations are, at best, demeaning and, at worst, damaging to health, if not accompanied by meaningful and up-to-date alternatives.

Predating the pandemic, publishers and education specialists have long suggested best practices for accessible visual media, including guidelines for tactile graphics~\cite{hasty_guidelines_2011} and for describing ``complex images'' in natural language~\cite{w3c_wai_2019, gould_effective_2008}.
While valuable, visualization authors have yet to broadly adopt these practices, for lack of experience with accessible media, if not a lack of attention and resources.
Contemporary visualization research has primarily attended to color vision deficiency~\cite{chaparro_applications_2017,nunez_optimizing_2018, oliveira_towards_2013}, and has only recently begun to engage with non-visual alternatives~\cite{choi_visualizing_2019, lundgard_sociotechnical_2019} and with accessibility broadly~\cite{kim_accessible_2021, wu_understanding_2021}.
Parallel to these efforts, computer science researchers have been grappling with the engineering problem of automatically generating chart captions~\cite{demir_summarizing_2012, obeid_chart--text_2020, qian_formative_2020}.
While well-intentioned, these methods usually neither consult existing accessibility guidelines, nor do they evaluate their results empirically with their intended readership.
As a result, it is difficult to know how useful (or not) the resultant captions are, or how effectively they improve access to meaningful information.

In this paper, we make a two-fold contribution.
First, we extend existing accessibility guidelines by introducing a conceptual model for categorizing and comparing the semantic content conveyed by natural language descriptions of visualizations.
Developed through a grounded theory analysis of {2,147} natural language sentences, authored by over 120 participants in an online study~\sectionlabel{section:3}, our model spans four levels of semantic content: enumerating visualization construction properties (e.g., marks and encodings); reporting statistical concepts and relations (e.g., extrema and correlations); identifying perceptual and cognitive phenomena (e.g., complex trends and patterns); and elucidating domain-specific insights (e.g., social and political context)~\sectionlabel{section:4}.
% (Moreover, we offer \emph{computational considerations} for generating natural language sentences at each level, such as access to the visualization specification, backing dataset, human perception, and/or domain-specific knowledge.)
Second, we demonstrate how this model can be applied to evaluate the effectiveness of visualization descriptions, by comparing different semantic content levels and reader groups.
We conduct a mixed-methods evaluation in which a group of {30} blind and {90} sighted readers rank the usefulness of descriptions authored at varying content levels~\sectionlabel{section:5}.
Analyzing the resultant {3,600} ranked descriptions, we find significant differences in the content favored by these reader groups: while both groups generally prefer mid-level semantic content, they sharply diverge in their rankings of both the lowest and highest levels of our model.

These findings, contextualized by readers' open-ended feedback, suggest that access to meaningful information is strongly {reader-specific}, and that captions for blind readers should aim to convey a chart's trends and statistics, rather than solely detailing its low-level design elements or high-level insights.
Our model of semantic content is not only \emph{descriptive} (categorizing what \emph{is} conveyed by visualizations) and \emph{evaluative} (helping us to study what \emph{should} be conveyed to whom) but also \emph{generative}~\cite{beaudouin-lafon_instrumental_2000, beaudouin-lafon_designing_2004}, pointing toward novel multimodal and accessible data representations~\sectionlabel{section:practicalImplications}.
Our work further opens a space of research on natural language as a data interface coequal with the language of graphics~\cite{bertin_semiology_1983}, calling back to the original linguistic and semiotic motivations at the heart of visualization theory and design~\sectionlabel{section:theoreticalImplications}.